% One producer, one consumer, one buffer, and the two ordering points.
\begin{tikzpicture}[font=\sffamily\small]
\node[hw, text width=26mm] (rdr) at (0,0) {USART3 RDR\\{\scriptsize one byte at a time}};
\node[fw, text width=28mm] (isr) at (3.6,0) {ring\_put\\{\scriptsize interrupt context}};
\node[kern, text width=28mm, minimum height=15mm] (buf) at (7.4,0)
  {ring\_t\\{\scriptsize 256 bytes\\head, tail, drops}};
\node[fw, text width=28mm] (get) at (11.2,0) {ring\_get\\{\scriptsize main loop}};

\node[ext, text width=26mm] (con) at (3.6,-2.9) {console\\{\scriptsize the drop count is printed here}};
\node[user, text width=28mm] (app) at (11.2,-2.9) {frame parser\\{\scriptsize chapter 5}};

\begin{scope}[on background layer]
\node[layer, fit=(isr)(buf)(get), inner ysep=8pt,
      label={[layerlabel, anchor=south west]north west:written in this chapter}] (dl) {};
\end{scope}

\draw[flow] (rdr) -- (isr);
\draw[flow] (isr) -- node[lbl, above]{write} (buf);
\draw[flow] (buf) -- node[lbl, above]{read} (get);
\draw[flow] (get) -- (app);
\draw[flow] (isr.south) -- node[lbl, right]{drops} (con.north);

\node[umlnote, text width=42mm, anchor=south] at (5.5,2.0)
  {Release. The data byte reaches memory before the head that publishes it.
   A volatile index alone does not give you this.};
\draw[dotted, draw=linecol] (5.5,1.95) -- (5.5,1.15);

\node[umlnote, text width=42mm, anchor=north] at (7.4,-2.1)
  {Acquire. The head is read before the slot it refers to. Both stale reads
   are conservative: never a read of an unwritten slot.};
\draw[dotted, draw=linecol] (7.4,-2.05) -- (7.4,-0.8);
\end{tikzpicture}
